\section{Conclusiones}
\label{sec:conclusiones}

\begin{table}[htbp]
    \centering
    \caption{Conclusiones del an\'{a}lisis de flexibilidad de la l\'{i}nea SIM-007.}
    \label{tab:conclusiones}
    \setcounter{itemcount}{0}
    \begin{tabularx}{\textwidth}{@{}NX@{}}
        \toprule
        \multicolumn{1}{c}{\textbf{Item}} & \textbf{Conclusion} \\
        \midrule
        & La l\'{i}nea de la bomba de vac\'{i}o nueva (SIM-007) cumple los criterios de esfuerzo de ASME B31.3-2016 en los nueve casos de carga evaluados: la evaluaci\'{o}n de cumplimiento de c\'{o}digo del modelo arroj\'{o} CODE COMPLIANCE EVALUATION PASSED. \\
        & El esfuerzo de c\'{o}digo m\'{a}ximo es \num{191915.1}~\si{kPa} en el nodo 100, caso 9 (EXP) L9=L1-L3, con una relaci\'{o}n de utilizaci\'{o}n de 70.1~\% frente al admisible de \num{273789.8}~\si{kPa}; la condici\'{o}n que gobierna el dise\~{n}o es el rango de expansi\'{o}n t\'{e}rmica entre las condiciones de operaci\'{o}n y la l\'{i}nea es la m\'{a}s exigida del paquete de flexibilidad del proyecto. \\
        & Los desplazamientos m\'{a}ximos del sistema son 3.5~\si{mm} (componente $DX$, caso 3 con $T_2$; hasta 3.4~\si{mm} en los casos de expansi\'{o}n) y las cargas en soportes y anclas son moderadas (carga vertical m\'{a}xima de operaci\'{o}n \num{1329}~\si{N} en el ancla del nodo 110), por lo que la configuraci\'{o}n de soporter\'{i}a del isom\'{e}trico es adecuada. \\
        & El margen disponible frente al admisible es de 29.9~\%: cualquier modificaci\'{o}n de la configuraci\'{o}n analizada reduce dicho margen y exige re-ejecutar el an\'{a}lisis. La comparaci\'{o}n de las cargas en la boquilla de la bomba con los admisibles del fabricante del equipo queda pendiente como verificaci\'{o}n externa, por no disponerse de dichos admisibles a la fecha. \\
        \bottomrule
    \end{tabularx}
\end{table}

En s\'{i}ntesis, la l\'{i}nea SIM-007 es t\'{e}cnicamente viable seg\'{u}n
ASME B31.3-2016, siendo la l\'{i}nea con mayor ratio de utilizaci\'{o}n del
paquete de flexibilidad, concentrado en el nodo 100. No se identifican riesgos
de sobre-esfuerzo, interferencias por expansi\'{o}n ni sobrecarga de soportes
en la configuraci\'{o}n analizada, por lo que no se requieren acciones
correctivas siempre que la construcci\'{o}n se apegue estrictamente al
isom\'{e}trico verificado.
